\documentclass{article}

\ifdefined\mathaicameraready
  \usepackage[dblblindworkshop, final]{neurips_2026}
\else
  \usepackage[dblblindworkshop]{neurips_2026}
\fi
\workshoptitle{The 6th Workshop on Mathematical Reasoning and AI}

\usepackage[utf8]{inputenc}
\usepackage[T1]{fontenc}
\usepackage{amsfonts}
\usepackage{amsmath}
\usepackage{amssymb}
\usepackage{colortbl}
\usepackage{amsthm}
\usepackage{booktabs}
\usepackage{fancyvrb}
\usepackage{etoolbox}
\usepackage{graphicx}
\usepackage{hyperref}
\usepackage{mathtools}
\usepackage{microtype}
\usepackage{tabularx}
\usepackage{url}
\usepackage{xcolor}
\usepackage{array}
\usepackage{float}
\usepackage{algorithm}
\usepackage{algorithmic}

% Consistent figure captions with the template's standard 7pt gap.
\AtBeginEnvironment{figure}{%
  \setlength{\abovecaptionskip}{7pt}%
  \setlength{\belowcaptionskip}{0pt}%
  \pretocmd{\caption}{\par\nointerlineskip}{}{}}
\hypersetup{hypertexnames=false,hidelinks,
  pdftitle={More than One Way to Skin a Cat: Preserving Verified Modes in RLVR},
  pdfauthor={Anonymous Authors}}
\newcolumntype{Y}{>{\raggedright\arraybackslash}X}
\newcolumntype{P}[1]{>{\raggedright\arraybackslash}p{#1}}
\newcommand{\1}{\mathbf 1}
\newcommand{\E}{\mathbb E}
\newcommand{\Bpos}{\mathcal B_x^+}
\newcommand{\mb}{\textsc{ModeBench}}
\definecolor{bettergreen}{RGB}{224,242,226}
\newcommand{\bettercell}[1]{\cellcolor{bettergreen}#1}
% Same five tints as each domain's card in Figure~\ref{fig:modebench-examples}
% (\path{ops/plot_paper_modebench_examples.py}, \texttt{DOMAIN_PANEL}), so a
% domain named in prose can be matched back to its card by colour alone.
\definecolor{domGraph}{HTML}{E8F1FA}
\definecolor{domCountdown}{HTML}{E7FBF6}
\definecolor{domPython}{HTML}{EFFBE7}
\definecolor{domMathIR}{HTML}{E7FBEE}
\definecolor{domPantry}{HTML}{E7ECFB}
\newcommand{\domlabel}[2]{\colorbox{#1}{\textbf{#2}}}
\newcommand{\xmode}{\mbox{legacy adaptive MaxEnt+ReplayDr.GRPO}}
\newcommand{\historicalt}{\mbox{historical multi-component treatment}}
\newcommand{\Historicalt}{\mbox{Historical multi-component treatment}}
\DeclareMathOperator{\clip}{clip}
\newtheorem{lemma}{Lemma}[section]
\newtheorem{theorem}[lemma]{Theorem}
\newtheorem{corollary}[lemma]{Corollary}
\theoremstyle{remark}
\newtheorem{remark}[lemma]{Remark}


\theoremstyle{definition}
\newtheorem{definition}[lemma]{Definition}
\theoremstyle{plain}


\newcommand{\qwenmark}{Qwen}



